\documentclass[11pt]{article}
\usepackage{amsmath,amssymb,amsthm,hyperref}
\newtheorem{theorem}{Theorem}
\title{A counterexample to the perpendicular perturbation conjecture for normed planes}
\author{}\date{}
\begin{document}
\maketitle
Fish, King, Miller, Palsson and Wahlenmayer conjectured that, for each
non-strictly convex norm on the plane, all sufficiently large line-like
crescent configurations are perpendicular perturbations
\cite[Conjecture 6.5]{fish}. Their Theorem 1.5 proves this for the maximum norm.
The construction below uses one norm for all sizes.
\begin{theorem}
Let $a=4/5$, $b=6/5$, and equip the plane with the norm
\[
 N(x,y)=\max\{(x^2+y^2)^{1/2},(ax^2+by^2)^{1/2},4|x+y|/5\}.
\]
For every $n\ge4$, put $h=\pi/(1000n)$ and
$P_j=(\cos\theta_j,\sin\theta_j)$ with
$\theta_j=\pi/2+jh$, $1\le j\le n$.
These points form a line-like crescent configuration in $N$, and no family
of parallel displacements takes them to equally spaced collinear points.
\end{theorem}
\begin{proof}
The maximum is a norm by the triangle inequality and positive definiteness of
its Euclidean term. The point $(5/8,5/8)$ lies strictly inside both quadratic
unit disks and on the supporting line $x+y=5/4$. Thus the unit circle of
$N$ contains a nondegenerate segment.

\paragraph{Distances.} Write $\eta=1/100$. All the angles lie in
$(\pi/2,\pi/2+\eta)$. For $i<j$, the direction of the chord $P_j-P_i$
has angle $(\theta_i+\theta_j)/2+\pi/2$. Its Euclidean unit vector $u$
satisfies $|u_y|<\eta$. Therefore $au_x^2+bu_y^2=a+(b-a)u_y^2<1$
and $4|u_x+u_y|/5<1$. Its $N$-length equals its Euclidean length.
In particular $N(P_j-P_i)=2\sin((j-i)h/2)$, strictly increasing in
the index gap. This gives precisely the equal-distance pattern of equally
spaced collinear points, as required by Definitions 2.8 and 2.10 of
\cite{fish}. No three points are collinear since they lie on a Euclidean circle.

\paragraph{Circumcircles.} Write $E(x,y)=(ax^2+by^2)^{1/2}$.
Any three chosen points have a unique $E$-circumcenter $c=(u,v)$
by transformation to the Euclidean metric. For its radius $R$ and
$z=e^{i\theta}$, the equation
\[
 a(\cos\theta-u)^2+b(\sin\theta-v)^2=R^2
\]
becomes, after multiplication by $z^2$, a quartic with leading and
constant coefficient $(a-b)/4$ and cubic coefficient $-au+ibv$.
If the three given roots are $z_1,z_2,z_3$, the remaining root is
$z_4=(z_1z_2z_3)^{-1}$. Putting $z_j=e^{i(\pi/2+\delta_j)}$
for $j\le3$ gives $z_4=e^{i(\pi/2-\sum_{j=1}^3\delta_j)}$.
Since $0<\delta_j<\eta$, this fourth root is outside the chosen arc.
Thus the $E$-circle passes through exactly three of the chosen points.

It remains to check circles of the actual norm $N$. Set $S=\sum_{j=1}^4z_j$.
Vieta's formula gives $au-ibv=(a-b)S/4$. The chord bound
$|e^{it}-1|\le |t|$ gives $|S-4i|<6\eta$. Hence
$|u|<0.0075$ and $|v-1/3|<0.005$. For \emph{every} chosen point let
$(X,Y)=P_j-c$. Then $|X|<0.0175$ and $Y>0.66$, so
$|X|<0.03Y$. Consequently
\[
 E(X,Y)^2-(X^2+Y^2)=(Y^2-X^2)/5>0,
 \]
\[
4|X+Y|/5\le 0.824Y < E(X,Y).
\]
Thus $N=E$ in neighborhoods of all these vectors. The $N$-circle with center
$c$ and radius $R$ contains the triple and excludes all other chosen points.
Moreover, in neighborhoods of the triple this circle is a strictly curved
ellipse. By Martini, Swanepoel and Wei\ss\ \cite[Proposition 21]{survey}, the
boundaries of two distinct positive homothets of a convex disk meet in the
union of at most two line segments, allowing singleton segments. At the
three strictly curved points no nondegenerate boundary segment is possible.
It follows that no other $N$-circle can contain the triple. This verifies
the no-four-on-a-norm-circle condition of \cite[Definition 4.4]{fish}.

\paragraph{Parallel displacements.} Suppose parallel displacements took the
points to equally spaced collinear points, allowing arbitrary relabeling.
A nonzero linear functional annihilating the displacements would project
the two sets to identical values. Its values on our set cannot all agree
by noncollinearity, so those values must be distinct elements of one
arithmetic progression of length $n$. For any linear functional,
$y_j=\ell(P_j)$ satisfies $y_{j+2}-2\cos(h)y_{j+1}+y_j=0$.
Subtracting at consecutive indices yields
\[
 (y_4-y_1)/(y_3-y_2)=1+2\cos h.
\]
The left side would be rational. But $2\cos h=e^{ih}+e^{-ih}$ is an
algebraic integer because $e^{ih}$ is a root of unity. It lies strictly
between 1 and 2, so is irrational. This excludes every perpendicular
perturbation allowed by \cite[Definition 4.1]{fish}.
\end{proof}
\paragraph{Relation to prior constructions.}
The circular-arc mechanism for the distance equalities occurs in
\cite[Theorem 1.3]{fish}. Their Example 4.3 has collinear triples for
$n\ge5$. Here the metric that determines the chord lengths and the
metric that determines the circumcircles of triples are active in different
directions. All circles excluded in the general-position argument are
circles of the one fixed norm $N$.
\begin{thebibliography}{9}
\bibitem{fish} S. Fish, D. King, S. J. Miller, E. A. Palsson and C. Wahlenmayer,
\emph{Crescent configurations in normed spaces}, Integers 20 (2020), A96.
\url{https://math.colgate.edu/~integers/u96/u96.pdf}
\bibitem{survey} H. Martini, K. J. Swanepoel and G. Wei\ss,
\emph{The geometry of Minkowski spaces---a survey. Part I},
Expositiones Mathematicae 19 (2001), 97--142.
\url{https://arxiv.org/abs/0708.2900}
\end{thebibliography}
\end{document}
